\begin{abstract}
Comparative claims about reinforcement-learning post-training of language
models are usually made between method labels: ``GRPO beats PPO'', ``DAPO beats
GRPO''. We report a set of negative and cautionary results showing that, in
the GRPO family, the label is not the unit of evidence: the training stack is
part of the result. A variance decomposition over 42 experiments attributes a
framework share ($\eta^2=0.546$) of the same order as the algorithm share
($0.558$). A matched PPO/GRPO contrast is a null (Welch $p=0.7605$, power
$0.06$) that reverses on a second model family. A same-stack audit of four
GRPO variants is inconclusive in every contrast. Under a fixed label, final
reward spans 5.0\% to 85.6\% across backends, and one label (``DAPO'') yields
opposite telemetry on two stacks; conversely, four labels on one fixed stack
land within a 0.034 band. We distil these exhibits into an eight-item
minimum reportable stack and a machine-readable registry whose diff tool flags
such label collisions from manifests alone.
\end{abstract}

\section{Introduction}

Group Relative Policy Optimization (GRPO) \citep{shao2024deepseekmath,guo2025deepseekr1}
replaces PPO's learned critic \citep{schulman2017proximal} with a baseline
computed from a group of $G$ sampled completions per prompt, and has spawned a
family of variants such as DAPO \citep{yu2025dapo}, Dr.\ GRPO
\citep{liu2025understanding} and GSPO \citep{qwen2025gspo}. Each is available
in several libraries, including TRL \citep{vonwerra2022trl}, veRL
\citep{sheng2024verl}, OpenRLHF \citep{hu2024openrlhf} and managed services such
as Tinker \citep{thinkingmachines2024tinker}, which differ in answer parser,
sampler, KL handling, clip range, LoRA targets \citep{hu2022lora}, precision
and evaluator. Deep-RL results are known to be sensitive to such details
\citep{henderson2018deep}, but papers in this area are not asked to report
them, so a gain attributed to a method label may belong to the stack.

This short paper asks a narrow question: \emph{how much of a reported GRPO-family
difference survives once the stack is held fixed?} Our answer, from our own
experiments, is: little that we could detect. We contribute (i) seven exhibits
of stack sensitivity, including three nulls that we report as nulls rather
than as equivalences, and (ii) an eight-item minimum reportable stack with an
auditor and a diff tool that operationalise the lesson.

\section{Setting and Protocol}

All runs use verifiable rewards, principally GSM8K exact match
\citep{cobbe2021training}, on open models from the Qwen and Llama families.
Where possible, each comparison fixes model, prompts, verifier, decoding
configuration, step budget and held-out evaluation, and varies one factor.
We use Welch or paired $t$-tests with Benjamini--Hochberg correction across the
test family, report effect sizes with 95\% confidence intervals, and state
observed power and the minimum detectable effect (MDE). An inconclusive test
is reported as inconclusive; it is not read as equivalence. Where we use the
zero-variance fraction (ZVF), it is the share of a batch's prompt groups in
which every completion received the same reward, so the group-relative
advantage is identically zero.

\section{Exhibits}

\begin{table*}[t]
\centering
\small
\begin{tabular}{@{}p{3.2cm}p{6.4cm}p{5.2cm}@{}}
\toprule
Exhibit & Result & Reading \\
\midrule
E-1 Variance decomposition (42 expts.) &
$\eta^2$: framework 0.546, algorithm 0.558, model family 0.471, scale 0.322 &
Framework term of the same order as the algorithm term; factors confounded \\
E-2 Matched PPO vs.\ GRPO (Qwen3-8B) &
Welch $p=0.7605$; $d=-0.14$, 95\% CI $[-1.02, 0.74]$; power 0.06; MDE 1.25 &
Failure to detect, not equivalence \\
E-3 Same contrast, Llama-3.1-8B-Instruct &
$d=12.75$, 95\% CI $[8.49, 17.00]$, $p<0.001$ &
Direction reverses across families; arms on different back-ends \\
E-4 Same-stack variant audit (8 paired seeds) &
DAPO $+0.00100$, GSPO $+0.00500$, Dr.\ GRPO $-0.00200$, AERO $-0.00075$ vs.\ GRPO; all inconclusive &
Exact MDE 0.0101 exceeds the 0.010 design margin \\
E-5 Same label, different backend &
Final training reward 5.0\% to 85.6\% ($17\times$) &
Backend bundled a different base checkpoint; not a backend-only effect \\
E-6 Same label ``DAPO'', two stacks &
Mean ZVF 0.00 (open trainer, true dynamic sampling) vs.\ 0.58 (closed stack, asymmetric-clip surrogate) &
One label, opposite telemetry \\
E-7 Four labels, one fixed stack (12 cells) &
Last-10 reward 0.744 DAPO, 0.742 Dr.\ GRPO, 0.723 GRPO, 0.710 GSPO; band 0.034 &
Comparable to seed variation; labels stop separating \\
\bottomrule
\end{tabular}
\caption{Stack-sensitivity exhibits. Rewards are GSM8K training or held-out
accuracy as stated in the text.}
\label{tab:exhibits}
\end{table*}

Table~\ref{tab:exhibits} summarises the exhibits.

\paragraph{The framework term is as large as the algorithm term (E-1).}
Across 42 experiments, a variance decomposition attributes
$\eta^2=0.546$ of outcome variation to library or framework, $0.558$ to
training algorithm, $0.471$ to model family and $0.322$ to model scale. These
factors are confounded with one another and with back-end, reward
implementation, hardware and task, so the shares are not additive causal
estimates. They do show that the framework is not a nuisance term that can be
ignored.

\paragraph{The headline comparison is a null (E-2, E-3).}
On Qwen3-8B, a matched PPO/GRPO contrast with ten runs per arm gives Welch
$p=0.7605$, the only one of twenty tests in our family that does not survive
Benjamini--Hochberg correction. The effect size is $d=-0.14$ with 95\% CI
$[-1.02, 0.74]$, observed power $0.06$ and an MDE of $1.25$ outcome units. The
interval admits effects in both directions; the study could rule out only a
very large difference. On Llama-3.1-8B-Instruct, the same contrast gives
$d=12.75$ ($[8.49, 17.00]$, $p<0.001$). The Qwen arms, however, ran on
different back-ends (managed service vs.\ an H100 cluster), so we treat both
rows as stack-conditioned evidence and withhold a directional algorithm
claim.

\paragraph{Same-stack variants do not separate (E-4).}
Holding a frozen stack fingerprint, we compared DAPO, GSPO, Dr.\ GRPO and AERO
against GRPO over eight paired seeds, evaluating each unit on 500 held-out
GSM8K items. All four contrasts are inconclusive, e.g.\ GSPO $+0.00500$
($[-0.00125, +0.01200]$, $p=0.210$) and AERO $-0.00075$
($[-0.00825, +0.00675]$, $p=0.858$). An earlier aggregator had concluded that
the differences ``disappear''; we retract that verdict because the locked
exact MDE ($0.0101$) exceeds the study's $0.010$ design margin. The
permitted conclusion is narrower: no arm produced a significant held-out
difference, and none met the pre-registered equivalence rule.

\paragraph{Labels hide stacks (E-5, E-6).}
Under a fixed label, changing the backend moved final training reward from
5.0\% to 85.6\%. We do not report this as a backend effect, because the same
inspection found that the managed backend had silently pinned a different base
checkpoint. That undisclosed bundling is itself the finding. Similarly, the
label ``DAPO'' produced mean ZVF 0.00 on an open trainer that implements true
dynamic sampling, and 0.58 on a closed stack that implements only an
asymmetric-clip surrogate.

\paragraph{Fixed stacks erase labels (E-7).}
The converse also holds. Running GRPO, DAPO, Dr.\ GRPO and GSPO on one
managed stack (Qwen3.5-4B, GSM8K, three seeds per arm), mean last-10 training
reward spans 0.710--0.744, a band of 0.034 that is comparable to seed-level
variation. On this corpus the algorithm term explains $\eta^2=0.0075$ of
reward variance, while the group-size axis explains far more.

\section{A Minimum Reportable Stack}

The exhibits motivate eight items that a stack-conditioned result must
report, each tied to a documented lever: (1)~the loss form (ratio level, clip
bounds, token mask, advantage normalisation, dynamic sampling), since E-6
shows one label with and without dynamic sampling; (2)~reference policy and
KL handling; (3)~sampler, backend and precision, including base-checkpoint
identity, motivated by E-5; (4)~the per-step ZVF and gradient-utilisation
trajectory, since a single mean conflates opposite regimes; (5)~the
group-size schedule; (6)~a held-out split disjoint from the reward
environment; (7)~a contamination check and parser-robustness probe; and
(8)~held-out pass@$k$ curves, reported with the evaluation rather than the run.
The items extend general reporting norms
\citep{pineau2020improving,mitchell2019modelcards,gebru2021datasheets} to the
specifics of group-relative RL.

We release the standard as a JSON schema, a registry of stack records and
variant-delta records, and a standard-library CLI. The schema separates an
unreported field (null) from a field reported as absent (e.g.\ ``no KL term''),
because only the latter is a checkable claim. An outcome-blind auditor scores
each record from 0 to 100 against the items; on the twenty seed records,
badges span 43--96. Toy-scale open runs score highest because every field is
reportable, not because they are better experiments. A diff tool maps a pair
of records onto a six-rung risk ladder, from identical manifests (R0) to
unauditable (R5, a differing item is null on one side). From manifests alone
it flags the E-6 ``DAPO'' pair as an R5 label-flip risk. In the registry,
the label ``DAPO'' is claimed by ten entries that cover four distinct
implementations.

\section{Discussion}

None of these exhibits shows that the GRPO variants are equivalent. They show
that, at the scale and power available to us, the comparisons that are
commonly reported could not have distinguished a method effect from a stack
effect. We think the productive response is procedural rather than
algorithmic: report the stack, audit it mechanically, and treat a label-only
comparison as unfalsifiable.

\section*{Limitations}

Our evidence comes from small-to-mid open models (roughly 0.5B--35B),
mathematics-style verifiable rewards and short training horizons (typically
30--50 steps on the managed service), so it says nothing about asymptotic
behaviour, preference-model rewards or frontier-scale training. Several
comparisons are single-seed or low-seed. The variance decomposition in E-1 is
confounded by design, and the managed service's GRPO internals are closed, so
results on it describe that platform's implementation rather than an abstract
algorithm. The 5.0\%--85.6\% exhibit mixes backend and checkpoint changes. The
registry's twenty seed records collapse to ten distinct fingerprints, and 20
of 23 sub-fields take at most two values, so the badge is coarse on this
corpus. The standard's eight items are motivated by our levers and have not
been validated by independent groups.

\section*{Ethics Statement}

This work uses public models and public mathematics benchmarks and involves
no human subjects or personal data. Its main risk is over-reading: our nulls
are underpowered and should not be cited as evidence that the methods are
equivalent.
